\begin{table}[!htbp]
\centering
\caption{Main model, 9th grade, Mathematics: ATT by relative period}
\label{tab:ap_event_base_9th_mat}
\vspace{-0.15cm}
\scriptsize
\setlength{\tabcolsep}{4pt}
\renewcommand{\arraystretch}{0.92}
\resizebox{\textwidth}{!}{%
\begin{tabular}{lcccc}
\toprule
Relative period & TWFE & CS21 & DR-DiD & S\&X approx. \\
\midrule
-4 & 0.072 (0.038) & 0.082 (0.044) & 0.060 (0.065) & 0.076 (0.053) \\
-3 & 0.066 (0.031) & 0.059 (0.032) & 0.037 (0.040) & 0.079 (0.038) \\
-2 & 0.010 (0.027) & 0.010 (0.027) & -0.003 (0.030) & 0.025 (0.032) \\
-1 & 0.000 & 0.000 & 0.000 & 0.000 \\
0 & 0.296 (0.034) & 0.290 (0.037) & 0.310 (0.038) & 0.311 (0.046) \\
1 & 0.422 (0.033) & 0.425 (0.036) & 0.442 (0.036) & 0.442 (0.050) \\
2 & 0.569 (0.035) & 0.583 (0.038) & 0.602 (0.039) & 0.595 (0.060) \\
3 & 0.716 (0.043) & 0.737 (0.048) & 0.751 (0.049) & 0.762 (0.080) \\
4 & 0.912 (0.052) & 0.890 (0.061) & 0.901 (0.062) & 0.907 (0.109) \\
\bottomrule
\end{tabular}
}
\begin{flushleft}
\footnotesize Note: each cell reports the ATT and standard error in parentheses, in standard deviations of annual proficiency. Period \(k=-1\) is normalized as the baseline and therefore appears without a standard error.
\end{flushleft}
\end{table}



